\documentclass[conference]{IEEEtran}
\IEEEoverridecommandlockouts
% The preceding line is only needed to identify funding in the first footnote. If that is unneeded, please comment it out.
\usepackage{cite}
\usepackage{amsmath,amssymb,amsfonts}
\usepackage{algorithmic}
\usepackage{graphicx}
\usepackage{textcomp}
\usepackage{xcolor}
\usepackage{booktabs}
\usepackage{microtype}
\usepackage[hyphens]{url}
\def\BibTeX{{\rm B\kern-.05em{\sc i\kern-.025em b}\kern-.08em
    T\kern-.1667em\lower.7ex\hbox{E}\kern-.125emX}}
\emergencystretch=1.5em
\Urlmuskip=0mu plus 1mu
\begin{document}

\title{Conference Paper Title\\}

\author{\IEEEauthorblockN{Khoa Le}
\IEEEauthorblockA{\textit{School of Computing} \\
\textit{University of Georgia}\\
Athens, USA \\
kl58277@uga.edu}
\and
\IEEEauthorblockN{Alexander Chen}
\IEEEauthorblockA{\textit{School of Computing} \\
\textit{University of Georgia}\\
Athens, USA \\
awc85026@uga.edu}
\and
\IEEEauthorblockN{Priyanka Shrestha}
\IEEEauthorblockA{\textit{School of Computing} \\
\textit{University of Georgia}\\
Athens, USA \\
ps21741@uga.edu}
}

\maketitle

\section{Introduction}
Multivariate time series forecasting has become increasingly important due to its wide range of applications in domains such as energy systems, finance, traffic prediction, and healthcare analytics. Despite recent advances in deep learning-based forecasting models, accurately capturing complex temporal dynamics and inter-channel relationships remains a significant challenge, particularly in non-stationary environments where data distributions shift over time, a phenomenon known as \textit{Temporal Distribution Shift} (TDS). 

Existing approaches handle inter-channel relationships through one of three strategies: Channel-Independent (CI), which treats each channel separately; Channel-Dependent (CD), which considers all channels jointly but is susceptible to noise; or Channel-Hard-Clustering (CHC) \cite{DGCformer}, which partitions channels into disjoint groups. Each of these has notable limitations in flexibility or robustness.

The DUET framework addresses these challenges through a dual-clustering mechanism operating on both the temporal and channel dimensions. Its Temporal Clustering Module (TCM) groups time series into latent distributions to model heterogeneous temporal patterns caused by TDS. In parallel, its Channel Clustering Module (CCM) captures inter-channel relationships in the frequency domain using a learnable Mahalanobis distance metric and sparsification, implementing a Channel-Soft-Clustering (CSC) strategy that is more flexible than prior hard-clustering approaches. These two modules are combined through a Fusion Module using masked multivariate attention.

In this work, we conduct an empirical analysis of DUET and evaluate it across nine benchmark datasets from the TFB suite, including ETTh1, ETTh2, ETTm1, ETTm2, Electricity, Exchange, Weather, Solar, and Traffic, as well as additional datasets including METR-LA, PEMS08, Wind, AQShunyi, AQWan, ZafNoo, and CzeLan. We further extend evaluation to the Sleep-EDF biomedical dataset to assess generalization to high-frequency signal domains. Additionally, we propose a lightweight local temporal residual branch based on depthwise Conv1D to capture short-term temporal dependencies and high frequency patterns that the original DUET architecture does not explicitly model.

\subsection{Objective}
The objective of this study is to evaluate and extend the DUET framework for multivariate time series forecasting. Specifically, we aim to 

(i) Replicate and evaluate DUET's performance on standard benchmark datasets from the TFB suite. 

(ii) investigate the impact of incorporating explicit local temporal modeling through a lightweight architectural modification.

(iii) Assess model generalization to biomedical time series by extending evaluation to the Sleep-EDF dataset, which contains multichannel EEG and EOG signals, and comparing against simple baselines to determine whether explicit local temporal modeling provides measurable benefit in high-frequency signal domains.

%\begin{figure}[htbp]
% \centerline{\includegraphics{duet_intro.png}}
% \caption{Example of a figure caption.}
% \label{fig}
% \end{figure}


\begin{figure}[htbp]
    \centering
    \includegraphics[width=0.75\columnwidth]{duet_intro.png}
    \caption{ The Temporal Clustering Module (TCM) and Channel Clustering Module (CCM) operate in parallel to capture heterogeneous temporal patterns and inter-channel dependencies, respectively. Their outputs are combined via a Fusion Module to produce the forecast. The local temporal branch (dashed) is our proposed extension, which adds a depthwise Conv1D residual path to capture short-term dependencies.}
    \label{fig:duet_diagram}
    \vspace{-0.5em}
\end{figure}

\subsection{Background}

Multivariate time series forecasting involves predicting future values of multiple correlated variables using historical observations. Compared to univariate forecasting, MTSF requires modeling both temporal dependencies and interactions among variables, making it inherently more complex.

One of the major challenges in MTSF is temporal non-stationarity, formalized as \textit{Temporal Distribution Shift} (TDS), where the statistical properties of the time series change over time. This leads to heterogeneous temporal patterns that a single global model struggles to capture. Another key challenge is the presence of complex inter-channel relationships, where variables influence each other in nonlinear and dynamic ways.

DUET addresses both challenges through a dual-clustering approach. It introduces a Temporal Clustering Module (TCM) that groups time series segments into latent distributions, enabling specialized pattern extractors to model heterogeneous temporal dynamics. In parallel, a Channel Clustering Module (CCM) captures inter-channel relationships in the frequency domain using a learnable distance metric and sparsification strategy. This implements a Channel-Soft-Clustering strategy where each channel attends only to channels that are beneficial for its prediction, mitigating the noise introduced by irrelevant channels.

The Fusion Module then combines the temporal features extracted by TCM with the channel mask from CCM using masked multivariate attention, and a final linear predictor maps the fused representation to the forecast horizon.

\subsection{Motivation}
While DUET effectively models global temporal structures and channel-wise dependencies, it does not explicitly incorporate an inductive bias for local temporal patterns. In many real-world scenarios, particularly in biomedical signals such as EEG, future values are strongly influenced by recent observations, short-term oscillations, and localized variations.

In such cases, models that rely primarily on global representations or clustering mechanisms may underfit fine-grained temporal dynamics. Convolution-based methods \cite{timesnet}, especially depthwise temporal convolutions, are well-suited for capturing short-term dependencies due to their localized receptive fields and computational efficiency.

Motivated by this limitation, we propose a lightweight extension to DUET by introducing a local temporal residual branch. This branch operates in parallel with the original architecture and captures high-frequency and short-range temporal patterns. By combining this local prediction with DUET’s output, the model can leverage both global and local temporal information.

Additionally, the original DUET paper focuses primarily on datasets from domains such as energy, traffic, and finance. To evaluate generalization, we extend the experiments to the Sleep-EDF dataset, where local temporal dynamics are more prominent. This allows us to assess whether explicit local modeling is beneficial in high-frequency and noisy signal environments.

\section{Literature Review}
Recent work in time series forecasting has explored various strategies to address temporal heterogeneity and cross-channel dependencies. Transformer-based models such as  iTransformer  \cite{informer} and Crossformer \cite{crossformer} utilize attention mechanisms to model long-range dependencies, often adopting either channel-independent or channel-dependent approaches. While effective, these models may be sensitive to noise or struggle with flexible channel interactions.

Other approaches focus on decomposition and efficiency. Models such as DLinear \cite{dlinear}, FITS \cite{FITS}, and TimeMixer \cite{timemixer} use linear transformations and decomposition techniques to capture temporal patterns with fewer parameters. However, they often assume relatively stable temporal structures and may not perform well under distribution shifts.

To address temporal non-stationarity, methods such as RevIN \cite{revin}, DAIN \cite{dain}, AdaRNN \cite{adaRNN}, and Non-stationary Transformers \cite{nstransformer} have been proposed. These approaches improve robustness by adapting to changing distributions, but typically handle temporal heterogeneity implicitly.

DUET builds upon these ideas by explicitly modeling both temporal distribution shifts and channel relationships through dual clustering. Its Temporal Clustering Module groups time series into latent distributions, while the Channel Clustering Module models inter-channel dependencies in the frequency domain using a learnable metric. This design enables DUET to achieve strong performance across diverse datasets.

In parallel, convolution-based architectures have been widely used to model local temporal dependencies. Their ability to capture short-term patterns complements attention-based and clustering-based methods. Our work integrates this ability into DUET by adding a local convolutional branch, resulting in a hybrid model that jointly captures global and local temporal dynamics.

\section{Experiments}
\subsection{Dataset Description}

\paragraph{Benchmark datasets}
We evaluate on multivariate long-term forecasting datasets from the DUET/TFB benchmark suite, including ETT variants (ETTh1/ETTh2/ETTm1/ETTm2), Exchange, Weather, AQ datasets, and others where scripts are provided.  
Each dataset contains synchronized multivariate observations over time, and forecasting is performed for horizons $H\in\{96,192,336,720\}$ when available.

\paragraph{Biomedical extension: Sleep-EDF.}
To test generalization beyond the original paper domain, we add Sleep-EDF \cite{sleep} as an external dataset.  
Raw PSG recordings are converted to DUET-compatible format (\texttt{date,value,cols}) using selected channels (EEG Fpz-Cz, EEG Pz-Oz, EOG horizontal), resampled to 1 Hz, and standardized per channel.  
Each subject-night is treated as one multivariate series.
\subsection{Experimental Setup}

\paragraph{Model and training protocol}
We use the official DUET implementation and rolling-forecast evaluation pipeline. For each run, the same train/validation/test policy is used as in the benchmark configuration, and we report normalized regression metrics from the generated report files.

\paragraph{Metrics.}
Primary metrics are \textbf{mse\_norm} and \textbf{mae\_norm}.  
Lower values indicate better forecasting performance.

\paragraph{Horizons and comparisons.}
We evaluate horizons $H\in\{96,192,336,720\}$.  
We report:
(i) DUET with our local temporal branch on benchmark datasets and 
(ii) Sleep-EDF results with all other datasets.
\paragraph{Architectural extension.}
We add an optional local residual branch to DUET:
depthwise temporal Conv1D + nonlinearity + projection from input length $L$ to horizon $H$, fused with DUET output using a learnable coefficient $\alpha$.  
The extension is controlled by \texttt{use\_local\_branch}, \texttt{local\_kernel\_size}, and \texttt{local\_alpha} in the main code.

\paragraph{Reproducibility}
All experiments are run through repository scripts under a fixed deterministic mode. Some large-scale runs (notably PEMS settings) encountered GPU Out-of-Memory, producing empty metric entries. Therefore, we omitted them for the final table. 
\subsection{Algorithm Description + Extension}

We use DUET, a multivariate time-series forecasting model that combines two ideas:
(1) temporal pattern clustering and (2) channel relationship modeling.
Given an input window $\mathbf{X}\in\mathbb{R}^{L\times C}$ (length $L$, channels $C$), DUET first applies a temporal pattern extractor with expert routing to model heterogeneous temporal dynamics. The resulting features are then processed by a channel-wise transformer block with a learned mask, so that dependencies among variables are captured explicitly. A linear prediction head maps latent features to the forecast horizon $H$. \textbf{Our extension}: the local temporal branch. To improve short-term pattern modeling, we add a lightweight residual branch in parallel to the main DUET path. This branch applies depthwise 1D temporal convolution per channel, followed by nonlinear activation and projection from $L$ to $H$. The final prediction is:
\[
\hat{\mathbf{Y}} = \hat{\mathbf{Y}}_{\text{DUET}} + \alpha \cdot \hat{\mathbf{Y}}_{\text{local}},
\]
where $\alpha$ is a learnable scalar (\texttt{local\_alpha}).  
This keeps DUET’s global/channel modeling intact while injecting local temporal inductive bias.

\subsection{Results}

\paragraph{Ablation of local branch.}
We compared DUET base vs. DUET+local branch (same data split/horizon).
On our Sleep-EDF and ETTh1 ablations, the local branch did not consistently improve metrics; in several runs the change was marginal.
This suggests the proposed branch is a valid architectural extension but requires additional tuning or dataset-specific adaptation to yield consistent gains.
\begin{table}[t]
\centering
\caption{Horizon-96 comparison between original DUET and DUET + Local Branch. Lower is better. Original DUET values are taken from Table 3 of the DUET paper}
\label{tab:h96_local_vs_original}
\footnotesize
\setlength{\tabcolsep}{3.5pt}
\renewcommand{\arraystretch}{0.95}
\begin{tabular}{lcc}
\toprule
Dataset & DUET (original) & DUET w. Local Branch (MSE/MAE) \\
\midrule
ETTh1 & \textbf{0.352} / \textbf{0.384} & 0.3614 / 0.3889 \\
ETTh2 & \textbf{0.270} / \textbf{0.336} & 0.2722 / 0.3390 \\
ETTm1 & \textbf{0.279} / \textbf{0.333} & 0.2898 / 0.3366 \\
ETTm2 & \textbf{0.161} / \textbf{0.248} & 0.1636 / 0.2539 \\
Electricity & 0.128 / 0.219 & \textbf{0.1274} / \textbf{0.2186} \\
Exchange & 0.080 / 0.198 & \textbf{0.0795} / \textbf{0.1975} \\
Solar & 0.169 / 0.195 & \textbf{0.1656} / \textbf{0.1943} \\
Traffic & \textbf{0.360} / \textbf{0.238} & 0.4417 / 0.2493 \\
Weather & \textbf{0.146} / \textbf{0.191} & 0.1492 / 0.1956 \\
SleepEDF & \textbf{1.2366} / \textbf{0.7979} & 1.2374 / 0.7981 \\
\bottomrule
\end{tabular}
\end{table}
\begin{table*}[t]
\centering
\caption{DUET + Local Branch results across benchmark datasets and Sleep-EDF extension. Lower is better. Traffic and some Sleep-EDF horizons are pending.}
\label{tab:duet_localbranch_with_sleep}
\footnotesize
\setlength{\tabcolsep}{3.5pt}
\renewcommand{\arraystretch}{0.95}
\resizebox{\textwidth}{!}{%
\begin{tabular}{lcccc}
\toprule
Dataset & Horizon 96 (MSE/MAE) & Horizon 192 (MSE/MAE) & Horizon 336 (MSE/MAE) & Horizon 720 (MSE/MAE) \\
\midrule
AQShunyi    & 0.7017 / 0.4692 & 0.7281 / 0.4853 & 0.7683 / 0.5040 & 0.8154 / 0.5225 \\
AQWan       & 0.7674 / 0.4505 & 0.8211 / 0.4723 & 0.8385 / 0.4790 & 0.9014 / 0.5104 \\
CzeLan      & 0.1621 / 0.1973 & 0.1944 / 0.2220 & 0.2129 / 0.2470 & 0.2582 / 0.2796 \\
ETTh1       & 0.3614 / 0.3889 & 0.4072 / 0.4146 & 0.4184 / 0.4328 & 0.5184 / 0.5095 \\
ETTh2       & 0.2722 / 0.3390 & 0.3340 / 0.3755 & 0.3677 / 0.4101 & 0.4576 / 0.4746 \\
ETTm1       & 0.2898 / 0.3366 & 0.3422 / 0.3648 & 0.3746 / 0.3885 & 0.4363 / 0.4182 \\
ETTm2       & 0.1636 / 0.2539 & 0.2157 / 0.2887 & 0.3001 / 0.3432 & 0.3580 / 0.3791 \\
Electricity & 0.1274 / 0.2186 & 0.1448 / 0.2350 & 0.1610 / 0.2525 & 0.1903 / 0.2785 \\
Exchange    & 0.0795 / 0.1975 & 0.1652 / 0.2909 & 0.2964 / 0.3941 & 0.5955 / 0.5855 \\
METR-LA     & 1.1277 / 0.5567 & 1.2780 / 0.6210 & 1.4302 / 0.6579 & 1.6655 / 0.7201 \\
PEMS08      & 0.1784 / 0.2086 & 0.2330 / 0.2021 & 0.3013 / 0.2187 & 1.4630 / 0.2960 \\
Solar       & 0.1656 / 0.1943 & 0.1926 / 0.2091 & 0.2090 / 0.2132 & 0.2269 / 0.2190 \\
Traffic     & 0.4417 / 0.2493 & 0.4753 / 0.2581 & 1.0819 / 0.4786 & 1.1147 / 0.4903 \\
Weather     & 0.1492 / 0.1956 & 0.1875 / 0.2295 & 0.2319 / 0.2647 & 0.3047 / 0.3215 \\
Wind        & 0.8745 / 0.6198 & 1.0437 / 0.7129 & 1.1496 / 0.7619 & 1.2318 / 0.8049 \\
ZafNoo      & 0.4298 / 0.3836 & 0.4786 / 0.4192 & 0.5293 / 0.4654 & 0.5691 / 0.4888 \\
\midrule
SleepEDF & 1.2374 / 0.7981 & 1.244 / 0.8004 & 1.2481 / 0.8021 & 1.2454 / 0.8025 \\
\bottomrule
\end{tabular}
}
\end{table*}

\paragraph{Discussion of failure cases.}
Some runs (PEMS variants datasets) produced empty metric fields due to GPU Out-of-Memory during channel-mask computation; these entries are marked as missing in our table, and we omitted them from the table.

\section{Conclusion and Discussion}
We added a local parallel residual branch to DUET to improve short-term pattern modeling while also adding a new biomedical dataset to test both the base model and our extension. The branch's output is multiplied by a learnable scalar and then combined with the base DUET’s prediction.

Although the addition of the parallel local branch did not disrupt DUET's dual-clustering mechanism, the ablation results show only minor changes for all datasets tested, including the Sleep-EDF dataset. The way we fuse the local and global branches may not be intricate enough to properly adapt to the datasets with the model potentially underweighting the local branch causing it to only change marginally.

We also had GPU Out-of-Memory (OOM) errors during the channel-mask computation phase when trying to run on the PEMS datasets. Computing pairwise inter-channel relationships scales quadratically with the number of channels making it sometimes uncompletable when there are too many channels like in PEMS. 

Future work should handle these two bottlenecks. To make the local branch effective the static $\alpha$ parameter may be changed to a dynamic mechanism that can weigh local versus global features. To make this architecture consistently runnable for datasets like PEMS, the channel masking step needs memory optimization or better hardware. Sparse attention approximations or chunked computations to keep VRAM usage bounded could potentially be used.

\bibliographystyle{IEEEtran}
\bibliography{references}

\end{document}